% GENERATED FILE. Edit canonical lessons, then run npm run generate:books.
% canonical-source-hash: fnv1a64:2b1fb9077adaa01a
% canonical-lessons: RU-C89-factory, RU-C89-office, RU-C89-corner, RU-C89-crossroads, RU-C89-traffic-lights

\chapter{Factory, Office, Corner, Crossroads, Traffic lights}
\label{ch:ru-factory-office-corner-crossroads-traffic-lights}
\hlchaptermodality{89}

\begin{chapteropening}
\textbf{By the end of this chapter:} \emph{I can say a factory, an office, a corner, a crossroads and traffic lights.}

Complete the last lesson of chapter 89: I can say a factory, an office, a corner, a crossroads and traffic lights.
\end{chapteropening}

\section[\texorpdfstring{\ru{завод} (zavód)}{завод (zavód)}]{\ru{завод} (zavód) — a factory}

\label{lesson:RU-C89-factory}



\begin{quote}
Before the new one: say the Russian for a swimming pool, then the Russian for a stadium.
\end{quote}

\subsection*{You'll want to know: \ru{завод}}
\textbf{\ru{завод}} (\emph{zavód}) — \textquotedblleft{}a factory\textquotedblright{}.

\begin{cousinweb}
\textbf{\ru{завод}} is \emph{\ru{за}-}, \textquotedblleft{}start\textquotedblright{}, with \emph{\ru{вод}}, the root of \emph{\ru{водить}}, \textquotedblleft{}to lead\textquotedblright{} (Vasmer gives \emph{\ru{водить}} with \emph{\ru{вести}}, \textquotedblleft{}to lead\textquotedblright{}): \emph{\ru{заводить}} is \textquotedblleft{}to set going\textquotedblright{}, and a \emph{\ru{завод}} is a works set going.
\end{cousinweb}

\begin{grammarlens}[title={What you've built}]
One more word for this chapter.
\end{grammarlens}

\subsection*{Your turn}
\begin{itemize}
\raggedright
  \item \emph{Say it:} \emph{zavód}
  \item \emph{Say it:} \emph{zavód}, once more
  \item \emph{Say it:} say \emph{zavód}
  \item \emph{Recall:} say the Russian for a swimming pool, then the Russian for a stadium, then say \emph{zavód} again
\end{itemize}

\begin{tcolorbox}[breakable,colback=teal!4,colframe=teal!35!black,title={Before you move on}]
What does \ru{завод} mean? (\textquotedblleft{}A factory\textquotedblright{}.) Say it once more.
\end{tcolorbox}
\section[\texorpdfstring{\ru{офис} (ófis)}{офис (ófis)}]{\ru{офис} (ófis) — an office}

\label{lesson:RU-C89-office}



\begin{quote}
Before the new one: say the Russian for a stadium, then the Russian for a factory.
\end{quote}

\subsection*{You'll want to know: \ru{офис}}
\textbf{\ru{офис}} (\emph{ófis}) — \textquotedblleft{}an office\textquotedblright{}.

\begin{cousinweb}
\textbf{\ru{офис}} comes from English \emph{office}, from Latin \emph{officium}, \textquotedblleft{}duty, service\textquotedblright{}.
\end{cousinweb}

\begin{grammarlens}[title={What you've built}]
One more word for this chapter.
\end{grammarlens}

\subsection*{Your turn}
\begin{itemize}
\raggedright
  \item \emph{Say it:} \emph{ófis}
  \item \emph{Say it:} \emph{ófis}, once more
  \item \emph{Say it:} say \emph{ófis}
  \item \emph{Recall:} say the Russian for a stadium, then the Russian for a factory, then say \emph{ófis} again
\end{itemize}

\begin{tcolorbox}[breakable,colback=teal!4,colframe=teal!35!black,title={Before you move on}]
What does \ru{офис} mean? (\textquotedblleft{}An office\textquotedblright{}.) Say it once more.
\end{tcolorbox}
\section[\texorpdfstring{\ru{угол} (úgol)}{угол (úgol)}]{\ru{угол} (úgol) — a corner}

\label{lesson:RU-C89-corner}



\begin{quote}
Before the new one: say the Russian for a factory, then the Russian for an office.
\end{quote}

\subsection*{You'll want to know: \ru{угол}}
\textbf{\ru{угол}} (\emph{úgol}) — \textquotedblleft{}a corner\textquotedblright{}.

\begin{cousinweb}
\textbf{\ru{угол}} $\leftarrow$ Proto-Slavic \emph{*ǫgŭlŭ} — inherited, a cousin of Latin \emph{angulus}, \textquotedblleft{}corner\textquotedblright{}, and Sanskrit \emph{aṅguli-}, \textquotedblleft{}finger\textquotedblright{} (Vasmer).
\end{cousinweb}

\begin{grammarlens}[title={What you've built}]
One more word for this chapter.
\end{grammarlens}

\subsection*{Your turn}
\begin{itemize}
\raggedright
  \item \emph{Say it:} \emph{úgol}
  \item \emph{Say it:} \emph{úgol}, once more
  \item \emph{Say it:} say \emph{úgol}
  \item \emph{Recall:} say the Russian for a factory, then the Russian for an office, then say \emph{úgol} again
\end{itemize}

\begin{tcolorbox}[breakable,colback=teal!4,colframe=teal!35!black,title={Before you move on}]
What does \ru{угол} mean? (\textquotedblleft{}A corner\textquotedblright{}.) Say it once more.
\end{tcolorbox}
\section[\texorpdfstring{\ru{перекрёсток} (perekryóstok)}{перекрёсток (perekryóstok)}]{\ru{перекрёсток} (perekryóstok) — a crossroads}

\label{lesson:RU-C89-crossroads}



\begin{quote}
Before the new one: say the Russian for an office, then the Russian for a corner.
\end{quote}

\subsection*{You'll want to know: \ru{перекрёсток}}
\textbf{\ru{перекрёсток}} (\emph{perekryóstok}) — \textquotedblleft{}a crossroads\textquotedblright{}.

\begin{cousinweb}
\textbf{\ru{перекрёсток}} is \emph{\ru{пере}-}, \textquotedblleft{}across\textquotedblright{}, with \emph{\ru{крест}}, \textquotedblleft{}a cross\textquotedblright{}, and \emph{-\ru{ок}}: where roads cross. Vasmer says \emph{\ru{крест}} first meant \textquotedblleft{}Christ\textquotedblright{}, from Old High German \emph{krist}, then \textquotedblleft{}crucifix\textquotedblright{}, then \textquotedblleft{}cross\textquotedblright{}.
\end{cousinweb}

\begin{grammarlens}[title={What you've built}]
One more word for this chapter.
\end{grammarlens}

\subsection*{Your turn}
\begin{itemize}
\raggedright
  \item \emph{Say it:} \emph{perekryóstok}
  \item \emph{Say it:} \emph{perekryóstok}, once more
  \item \emph{Say it:} say \emph{perekryóstok}
  \item \emph{Recall:} say the Russian for an office, then the Russian for a corner, then say \emph{perekryóstok} again
\end{itemize}

\begin{tcolorbox}[breakable,colback=teal!4,colframe=teal!35!black,title={Before you move on}]
What does \ru{перекрёсток} mean? (\textquotedblleft{}A crossroads\textquotedblright{}.) Say it once more.
\end{tcolorbox}
\section[\texorpdfstring{\ru{светофор} (svetofór)}{светофор (svetofór)}]{\ru{светофор} (svetofór) — traffic lights}

\label{lesson:RU-C89-traffic-lights}



\begin{quote}
Before the new one: say the Russian for a corner, then the Russian for a crossroads.
\end{quote}

\subsection*{You'll want to know: \ru{светофор}}
\textbf{\ru{светофор}} (\emph{svetofór}) — \textquotedblleft{}traffic lights\textquotedblright{}.

\begin{cousinweb}
\textbf{\ru{светофор}} is \emph{\ru{свет}}, \textquotedblleft{}light\textquotedblright{}, with \emph{-\ru{фор}}, from Greek \emph{-phóros}, \textquotedblleft{}carrying\textquotedblright{}: \textquotedblleft{}light-bearer\textquotedblright{}, a Soviet coinage.
\end{cousinweb}

\begin{grammarlens}[title={What you've built}]
One more word for this chapter.
\end{grammarlens}

\subsection*{Your turn}
\begin{itemize}
\raggedright
  \item \emph{Say it:} \emph{svetofór}
  \item \emph{Say it:} \emph{svetofór}, once more
  \item \emph{Say it:} say \emph{svetofór}
  \item \emph{Recall:} say the Russian for a corner, then the Russian for a crossroads, then say \emph{svetofór} again
\end{itemize}

\begin{tcolorbox}[breakable,colback=teal!4,colframe=teal!35!black,title={Before you move on}]
What does \ru{светофор} mean? (\textquotedblleft{}Traffic lights\textquotedblright{}.) Say it once more.
\end{tcolorbox}
